%! Confidence Intervals for the Difference of Two Means — Unit 7, Lesson 7.6 — Warm-Up (Answer Key)
\documentclass[10pt]{article}
\usepackage{apstats-article}
\usepackage{apstats-key}   % pulls in -boxes; do NOT also load -boxes
\newcommand{\TallMath}[1]{$\displaystyle #1\rule[-1.4em]{0pt}{3.2em}$}

\begin{document}

\pageheader{Unit 7, Lesson 7.6}{Warm-Up --- Answer Key}
\namedateperiod

\vspace{0.15in}

\begin{notesbox}{Spiral Review --- One-Sample $t$-Interval (Lessons 7.2--7.3)}

A school dietitian claims that students consume an average of 850 calories at lunch. A
random sample of 22 students recorded their lunch calories. The results were
$\bar x = 920$ cal and $s = 130$ cal. A dotplot shows no strong skew or outliers.

\textbf{1.} Name the inference method and give the degrees of freedom.

\quad Method: \ans{one-sample $t$-interval for $\mu$} \quad $df = $ \ans{21}

\textbf{2.} Verify both conditions.
\begin{itemize}
  \item Random: \ansline{Random sample of 22 students --- met.}
  \item Normal: \ansline{$n = 22 < 30$; dotplot shows no strong skew or outliers --- met.}
\end{itemize}

\textbf{3.} Compute the standard error and margin of error for a 95\% confidence interval
($t^* = 2.080$). Show your work.

\vspace{0.1in}
$SE = \dfrac{s}{\sqrt{n}} = \dfrac{{\color{keyred}\mathbf{130}}}{\sqrt{{\color{keyred}\mathbf{22}}}} \approx $ \ans{27.72}
\hfill
$ME = t^* \times SE = {\color{keyred}\mathbf{2.080}} \times {\color{keyred}\mathbf{27.72}} \approx $ \ans{57.66}

\textbf{4.} Write the 95\% confidence interval and interpret it in context.

\quad CI: $\bar x \pm ME = {\color{keyred}\mathbf{920}} \pm {\color{keyred}\mathbf{57.66}} = ($ \ans{862} , \ans{978} $)$

\quad Interpretation: \ansline{We are 95\% confident the true mean lunch calorie intake for students is between 862 and 978 calories.}

\end{notesbox}

\end{document}
